% WS3 report fragment (issue #39): candidate verification, big-bucket handling,
% SON frequent-tagset blocking, E4/E6/E7 results and limitations, and WS3 contribution.
% No preamble. Included via \input in docs/report/main.tex (WS3).
% Defines \wsthreeMethodology (Section 5 part), \wsthreeResults (Section 8 part),
% \wsthreeAnalysis (Section 9 part), and \wsthreeContribution (Section 11 part).
% Every empirical number carries a "% source:" comment; "% TODO" marks pending cluster runs.
%
% Citation keys in docs/report/refs.bib:
%   \cite{salton1988tfidf, savasere1995son, agrawal1994apriori}

\newcommand{\wsthreeMethodology}{%
\subsection{Candidate Verification and Re-Ranking (E4)}
LSH candidate pairs (Contract C4) carry only a noisy MinHash estimate $\hat{J}(A, B)$ with standard
deviation $\sqrt{J(1-J)/n}$ and positive selection bias toward pairs whose random band hashes
collide. In Stage C4 $\to$ C5 (\texttt{src/verification/verify.py}), we join candidate IDs against
the C2 shingle sets and compute exact Jaccard similarity:
\begin{equation}
  J(A, B) = \frac{|A \cap B|}{|A \cup B|}.
\end{equation}
A shared \texttt{Ranker} protocol (\texttt{src/verification/rankers.py}) supports both
\texttt{EstimatedRanker} and \texttt{ExactJaccardRanker}. As a stretch extension (\#40),
\texttt{TfidfRanker} (\texttt{src/verification/rerank.py}) computes unit-normalized TF-IDF
cosine similarity \cite{salton1988tfidf} over the top-$k$ verified candidates per query to resolve
lexical ties using rare-term weights.

\subsection{Big-Bucket Skew Mitigation (E7)}
Popular or boilerplate title shingles can push many questions into a single LSH bucket, turning a
single Spark partition into a quadratic bottleneck. Through Contract Hook 1
(\texttt{--bucket-policy}), \texttt{src/verification/buckets.py} profiles bucket sizes (max size,
total implied pairs, and largest-bucket pair share) and compares three policies without modifying
\texttt{lsh.py}:
\begin{enumerate}
  \item \textbf{Pass-through (\texttt{none}):} keep all buckets as an unmitigated baseline.
  \item \textbf{Size limit (\texttt{size\_limit}):} drop any bucket with more than $L$ members.
  \item \textbf{Drop common shingles (\texttt{drop\_common\_shingles}):} prune shingles whose
        document frequency exceeds \texttt{drop\_common\_share} in \texttt{ShingleConfig}.
\end{enumerate}
As a stretch alternative (\#41), \texttt{salt\_split} was planned to partition oversized buckets
into $k$ chunks and emit chunk-pair sub-buckets so all candidate pairs are preserved while work
spreads across tasks, but it was cut at the checkpoint to prioritize core E6 and E7 cluster runs.

\subsection{Frequent-Tagset Blocking with SON (E6)}
To study whether topic partitioning reduces LSH work (Q4), we implement the two-pass Savasere,
Omiecinski, and Navathe (SON) algorithm \cite{savasere1995son} with Apriori
\cite{agrawal1994apriori} in Spark RDDs (\texttt{src/blocking/son.py}). Pass~1 mines locally
frequent tagsets per partition at scaled threshold $s \cdot |P_i|$; Pass~2 counts exact global
support across all baskets. Using the discovered frequent tagsets, \texttt{src/blocking/\allowbreak blocking.py}
assigns each question to its largest frequent tagset (breaking ties by lowest global support, or
\texttt{other} if none match) and outputs Contract C9 (\texttt{post\_id}, \texttt{group}) for
\texttt{lsh.py --groups}, alongside \texttt{single\_tag} and \texttt{none} baselines.
}

\newcommand{\wsthreeResults}{%
\subsection{Verification, Blocking, and Bucket Results (E4, E6, E7)}
On the 5,000-question Super User sample (word 1-gram title shingles, $n=128$, $b=32$, $r=4$), exact
Jaccard verification re-scores 13,417 top-10 candidate pairs across 2,780 queries with candidates
% source: results/e4_verification/ranking_comparison.csv
(Table~\ref{tab:e4-verification}). Estimated MinHash overstates similarity on average ($0.265423$ vs.
$0.250672$ exact), and exact verification changes the top-1 candidate for $17.70\%$ of queries
% source: results/e4_verification/ranking_comparison.csv
(top-1 agreement $82.30\%$, mean absolute rank shift $0.937854$ positions).

\begin{table}[h]
  \centering
  \small
  \setlength{\tabcolsep}{3pt}
  \caption{E4 ranking comparison on 5,000 Super User questions ($2{,}780$ queries, $13{,}417$ ranked pairs).}
  \label{tab:e4-verification}
  \begin{tabular}{@{}lccc@{}}
    \toprule
    Method & Mean score & Top-1 agree. & Rank shift \\
    \midrule
    Estimated & $0.265423$ & $0.823022$ & $0.937854$ \\ % source: results/e4_verification/ranking_comparison.csv
    Exact Jaccard     & $0.250672$ & $1.000000$ & $0.000000$ \\ % source: results/e4_verification/ranking_comparison.csv
    \bottomrule
  \end{tabular}
\end{table}

For E6 SON mining on Super User ($516{,}669$ questions),
% source: docs/report/ws1_intro_data_eval.tex (Super User question count)
two-pass SON finds $52$ frequent tagsets in $77.4713$\,s versus $34.5360$\,s for MLlib
\texttt{FPGrowth}
% source: results/e6_blocking/frequent_tagsets.csv
while matching its frequent itemsets exactly. On the full Stack Overflow \texttt{python} subset
($2{,}204{,}200$ questions, $102{,}878$ answer-key pairs, $n=128$, $b=32$, $r=4$), SON mining at
support $0.005$ (\texttt{max\_size=3}) discovers $141$ frequent tagsets in $64.6821$\,s versus
$9.8758$\,s for \texttt{FPGrowth} and assigns all questions across $72$ groups ($0.00\%$ in
\texttt{other}, $33.01\%$ in \texttt{python}).
% source: results/e6_blocking/frequent_tagsets_stackoverflow.csv, results/e6_blocking/README.md
Table~\ref{tab:e6-blocking} compares blocking modes: \texttt{single\_tag} (most frequent tag
other than \texttt{python}) cuts candidates by $68.4\%$ and runtime by $22.9\%$ at a
\texttt{recall\_cost} of $0.620891$, while \texttt{son} cuts candidates by $60.0\%$ at
\texttt{recall\_cost} $0.544072$.

\begin{table}[h]
  \centering
  \small
  \setlength{\tabcolsep}{3pt}
  \caption{E6 tag-blocking comparison on $2{,}204{,}200$ Stack Overflow \texttt{python} questions ($n=128$, $b=32$, $r=4$).}
  \label{tab:e6-blocking}
  \begin{tabular}{@{}lrrrrrr@{}}
    \toprule
    Mode & Time (s) & Shuffle (GB) & Candidates & Recall@10 & MAP & Recall cost \\
    \midrule
    \texttt{none} & $3126.98$ & $174.44$ & $740{,}022{,}352$ & $0.037762$ & $0.023023$ & $0.000000$ \\ % source: results/e6_blocking/blocking_comparison.csv
    \texttt{single\_tag} & $2412.46$ & $165.14$ & $233{,}502{,}916$ & $0.023807$ & $0.015507$ & $0.620891$ \\ % source: results/e6_blocking/blocking_comparison.csv
    \texttt{son} & $2880.62$ & $168.62$ & $296{,}065{,}255$ & $0.027072$ & $0.017176$ & $0.544072$ \\ % source: results/e6_blocking/blocking_comparison.csv
    \bottomrule
  \end{tabular}
\end{table}

For E7 on the full Stack Overflow \texttt{python} subset (Table~\ref{tab:e7-buckets}),
\texttt{drop\_common\_shingles} (\texttt{drop\_common\_share=0.05}) is $25.5\%$ faster overall
($2261.46$\,s vs. $3033.79$\,s), cuts candidate pairs by $21.9\%$ ($578{,}143{,}060$ vs.
$740{,}022{,}352$), and raises Recall@10 by $+25.8\%$ ($0.047498$ vs. $0.037762$), whereas
\texttt{size\_limit} ($1000$) bounds worst-case task skew ($148.39$\,s slowest task vs.
$403.50$\,s).
% source: results/e7_buckets/bucket_comparison.csv, results/e7_buckets/README.md

\begin{table}[h]
  \centering
  \small
  \setlength{\tabcolsep}{3pt}
  \caption{E7 big-bucket strategy comparison on $2{,}204{,}200$ Stack Overflow \texttt{python} questions.}
  \label{tab:e7-buckets}
  \begin{tabular}{@{}lrrrrrr@{}}
    \toprule
    Strategy & Time (s) & Slowest (s) & Median (s) & Shuffle (GB) & Candidates & Recall@10 \\
    \midrule
    \texttt{size\_limit(1000)} & $3033.79$ & $148.39$ & $6.26$ & $141.18$ & $740{,}022{,}352$ & $0.037762$ \\ % source: results/e7_buckets/bucket_comparison.csv
    \texttt{drop\_common\_shingles(0.05)} & $2261.46$ & $403.50$ & $1.55$ & $151.13$ & $578{,}143{,}060$ & $0.047498$ \\ % source: results/e7_buckets/bucket_comparison.csv
    \bottomrule
  \end{tabular}
\end{table}
}

\newcommand{\wsthreeAnalysis}{%
\subsection{Verification, Blocking, and Skew Challenges}
\paragraph{Why verification matters.} LSH candidate generation admits pairs whose band hashes happen
to collide, inflating \texttt{est\_sim} above true Jaccard. Exact C2 verification eliminates this
variance cheaply after LSH has already pruned the quadratic search space.

\paragraph{Big-bucket skew and checkpoint cuts.} Skewed buckets degrade Spark stage tail latency
and executor memory. \texttt{size\_limit} bounds worst-case bucket work at the cost of dropping
oversized buckets, whereas \texttt{drop\_common\_shingles} prevents common tokens from forming huge
buckets earlier in Stage C1 $\to$ C2. For E8 scaling runs, \texttt{drop\_common\_shingles} is the
primary recall-preserving strategy, paired with \texttt{size\_limit} as a tail-skew guard. At the
checkpoint, streaming, AI-impact experiments, and E7 \texttt{salt\_split} (\#41) were cut so
cluster budget focused on core E6 and E7 runs, while E4 TF-IDF re-ranking (\#40) was completed.

\paragraph{Tag blocking trade-off.} Blocking by frequent tagsets shrinks candidate comparisons
within topics, but any duplicate pair whose questions do not share the same assigned group incurs
\texttt{recall\_cost} across group boundaries.
}

\newcommand{\wsthreeContribution}{%
\paragraph{Navjot Singh Shienh (WS3: Candidate Refinement, Blocking, and Report Assembly).}
Implemented exact Jaccard verification (\texttt{verify.py}, \texttt{rankers.py}, E4) and TF-IDF
cosine re-ranking (\texttt{rerank.py}, \#40); two-pass Spark RDD SON mining and tag-blocking
assignment (\texttt{son.py}, \texttt{blocking.py}, E6); big-bucket profiling and
\texttt{size\_limit} vs. \texttt{drop\_common\_shingles} comparison (\texttt{buckets.py}, E7);
Laptop~B Spark cluster integration; and the LaTeX report assembly and validation tooling
(\texttt{main.tex}, \texttt{refs.bib}, \texttt{check\_report.py}).
}
